\section{Introduction}\label{intro:section}

A geometric Hecke eigensheaf realizes a local system on a curve through modifications of bundles. At a puncture, the monodromy of the local system must be reflected in the level structure on the bundle stack. We construct such eigensheaves for $\SL_n$ with Borel level at several punctures and arbitrary unipotent tame monodromy. The sheaves are locally constructible and perverse, and their eigenisomorphisms extend coherently over collisions of the modification points.

\subsection{The theorem}
Let $n\ge2$, let $k=\overline{\mathbb F}_q$ have characteristic $p>n$, and fix a prime $\ell\ne p$. Put $E=\overline{\mathbb Q}_\ell$. Let $X_0/\mathbb F_q$ be a smooth projective geometrically connected curve of genus $g\ge2$, with distinct rational points $x_1,\ldots,x_s$, where $s\ge2$. Write
\[
 X=(X_0)_k,\qquad D=x_1+\cdots+x_s,\qquad U=X\setminus|D|.
\]
Set $G=\SL_n$ and $\Gd=\PGL_{n,E}$, and let $B\subset G$ be the upper-triangular Borel subgroup. All bundle stacks below refer to the split group over the geometric base field in use.

A parameter is a continuous homomorphism
\begin{equation}\label{intro:parameter}
 \rho:\pi_1^{\et}(U,\bar u)\longrightarrow\PGL_n(L),
\end{equation}
where $L/\mathbb Q_\ell$ is finite inside $E$. We assume that its image is Zariski dense, that it kills wild inertia at every $x_i$, and that
\begin{equation}\label{intro:monodromy}
 \rho(\gamma)=\exp\bigl(t_{\ell,i}(\gamma)N_i\bigr)
 \qquad(\gamma\in I_{x_i}).
\end{equation}
Here $t_{\ell,i}:I_{x_i}\to\mathbb Z_\ell$ is a chosen tame quotient and $N_i\in\mathfrak{pgl}_n(L)$ is nilpotent. The exponential is the finite nilpotent exponential. No regularity is assumed for $N_i$; it may be zero. For $V\in\Rep_E(\Gd)$, write $V_\rho$ for the corresponding lisse $E$-sheaf on $U$.

Let
\[
 \cA=\Bun_{\SL_n,B,D}(X)
\]
be the stack of $\SL_n$-bundles $P$ together with Borel reductions $\beta_i$ at the marked points. A geometric $E$-adic complex $M$ on $\cA$ is \emph{locally constructible perverse} if $f^*M[d]$ is bounded constructible and perverse for every smooth finite-type chart $f:S\to\cA$ of constant relative dimension $d$. This is a local condition on the stack.

The trace pairing identifies its cotangent vectors with Higgs fields
\begin{equation}\label{intro:cotangent}
 \phi\in H^0\bigl(X,\ad(P)\otimes\omega_X(D)\bigr),
 \qquad \Res_{x_i}(\phi)\in\Lie R_u(B_{\beta_i})
 \quad(1\le i\le s).
\end{equation}
Let $\Lambda_D$ be the cone of these fields that are nilpotent at the generic point of $X$. A singular-support containment in this cone means its smooth cotangent pullback on every smooth chart.

For a finite set $I$ and representations $V_i\in\Rep_E(\Gd)$, the functor $\Hecke_{I,(V_i)}$ modifies bundles at points of $U^I$. We use relative Satake normalization: on a Schubert cell of dimension $d_\lambda$ the kernel is $E[d_\lambda](d_\lambda/2)$, and there is no additional moving-leg shift. The conventions and collision maps are specified in Section~\ref{found:section}.

\Needspace{12\baselineskip}
\begin{theorem}\label{intro:main}
For every parameter \eqref{intro:parameter} satisfying the preceding hypotheses, there is a nonzero locally constructible perverse geometric $E$-adic sheaf $M$ on $\cA$ with
\[
 \SS(M)\subset\Lambda_D.
\]
For every finite set $I$ and every family $(V_i)_{i\in I}$, it has isomorphisms
\begin{equation}\label{intro:eigen}
 \Hecke_{I,(V_i)}(M)
 \simeq M\boxtimes\Bigl(\mathop{\boxtimes}_{i\in I}(V_i)_\rho\Bigr)
 \qquad\text{on }\cA\times U^I,
\end{equation}
natural in the representations and coherent for the tensor unit, convolution, permutations, and fusion on every collision diagonal. These isomorphisms use the relative normalization above.
\end{theorem}

The parameter and eigensheaf are geometric: the statement does not require a Frobenius structure. Local constructibility allows the cohomological bounds to depend on the finite-type chart.

\subsection{Context and antecedents}
The eigensheaf problem originates in Drinfeld's rank-two construction and
Laumon's geometric formulation for general linear groups
\cite{Drinfeld,Laumon}. Frenkel, Gaitsgory, and Vilonen reduced unramified
$\GL_n$ eigensheaf existence to a vanishing conjecture \cite{FGV}, which
Gaitsgory subsequently proved \cite{GaitsgoryVanishing}. Geometric Satake
identifies the local Hecke kernels with representations of the dual group
and supplies their tensor and fusion structure \cite{MV}. These results
established the geometric form of the passage from a prescribed local
system to an automorphic sheaf.

Our immediate antecedent is \cite[Theorem~1.1]{CT}, whose main theorem
has one marked point and regular-unipotent monodromy for a broader class
of groups. Here the group is $\SL_n$ with $p>n$, there are several marked
points, and every nilpotent boundary logarithm is allowed. The
two-specialization construction and the local arguments follow that
companion; we prove the simultaneous level geometry and compatibility
arguments needed for this extension.

For semisimple groups over $\mathbb C$, Beilinson and Drinfeld constructed
eigensheaves from opers by quantizing the Hitchin system \cite[\S0.2]{BD}.
The recent five-paper collaboration proves categorical forms of the
unramified correspondence in characteristic zero \cite{GLCI,GLCV}.
Arinkin, Gaitsgory, Kazhdan, Raskin, Rozenblyum, and Varshavsky introduced
the stack of local systems with restricted variation and its action on
sheaves with nilpotent singular support \cite{AGKRRV}.
Gaitsgory and Raskin use this framework to formulate the geometric adic
correspondence and compare characteristic zero with positive characteristic
\cite{GR}. The input needed here is the characteristic-zero restricted
equivalence \cite[Main Theorem~1.3.9(ii)]{GR}, through the unramified
eigencomplex construction of \cite[Lemma~8.1]{CT}. Thus the categorical
theory enters on a smooth projective curve before level structures are
introduced.

Nilpotent singular support links this input to ramification. The geometry
of the global nilpotent cone was developed by Laumon, Faltings, and Ginzburg
\cite{LaumonNilpotent,Faltings,GinzburgNilpotent}. Beilinson's theory of
singular support for constructible \'etale sheaves, and Barrett's extension
to adic coefficients, give the function tests used below
\cite{Beilinson,Barrett}. The Hecke detection argument in
\cite[\S\S20.5--20.8, Appendix~H]{AGKRRV} uses a cocharacter in the center
of a Levi subgroup and an isolated cotangent intersection. The field and
specialization arguments of \cite[\S\S3--4]{CT} adapt that geometry. We establish
the versions required for all the level stacks occurring here.
Geometric nearby cycles and their perverse exactness rest on
Hansen--Scholze and Gaitsgory--Raskin \cite{HS,GR}.

The characteristic-zero level arguments use Betti sheaf theory.
Nadler and Yun prove local constancy of moving Hecke operators on the
nilpotent category and construct its spectral action, already allowing
finitely many level points \cite[Theorems~6.2.2 and~6.3.7]{NY}.
Their automorphic gluing construction relates twisted nodal curves to
smoothings \cite{NYGluing}; its description of bundle types and branch
gluing is the geometric antecedent of our degeneration. Our proof also
needs a separate nonvanishing statement for the particular nearby cycles
being taken. At a puncture, the local spectral input comes from
Gaitsgory's central sheaves \cite{GaitsgoryCentral}, the affine-flag
equivalence of Arkhipov and Bezrukavnikov \cite{AB}, and Dhillon and
Taylor's extension to the universal monodromic setting \cite{DT}.
The last supplies the local kernels used in the descent argument of
\cite[\S7]{CT}, which we globalize with all the other levels retained.
Dhillon and Taylor's sequel proves the tame local Betti correspondence
\cite{DTLocal}.

Several-point ramification has earlier geometric realizations. Uit de Bos
gives an explicit $\GL_2$ correspondence for the projective line with four
punctures and pure irreducible parameters with unipotent monodromy
\cite{UitDeBos}. Shen proves a Hecke-compatible equivalence over an open
spectral locus for parabolic $\GL_n$ in positive characteristic, using
crystalline $D$-modules \cite{Shen}. In the characteristic-zero de Rham
setting, F{\ae}rgeman gives a construction for irreducible regular-singular
parameters conditional on the forthcoming factorization input and local
compatibility conjecture specified in that work
\cite[Theorem~1.4.1.1, \S\S1.3.7, 1.4.3, Conjecture~4.2.1]{Faergeman}.
In the cited version,
regular holonomicity and generic perversity in the ramified setting are
stated as expectations \cite[Remark~1.4.1.2]{Faergeman}. The present
theorem concerns geometric adic sheaves in positive characteristic and
includes local constructibility and perversity.


\subsection{Proof strategy and the several-point extensions}
The proof has two specializations. First we construct an eigensheaf with the prescribed level in characteristic zero. We then specialize it to the given curve in characteristic $p$. The intermediate sheaves must remain locally constructible, retain the full Hecke system, and remain nonzero. The last condition requires a separate argument: nearby cycles of a nonzero complex can vanish.

The first stage starts with any marked complex curve and a dense parameter with unipotent boundary monodromy. Take two copies of that curve and join corresponding marked points to form a nodal curve. A smoothing replaces all its nodes at once. On the second component, reverse the orientation of the underlying marked surface and transport the parameter. This makes the two boundary monodromies inverse at every pair of branches. Finite quotients of the parameter then glue at all nodes after introducing suitable cyclic stabilizers there. Their compatible lifts produce a dense unramified parameter on the smooth geometric generic curve, where the characteristic-zero existence theorem applies.

The cyclic stabilizer type of the bundles has a different role from the boundary monodromy of the parameter. Choosing that type inside an alcove turns the special bundle stack into a quotient of two stacks with \emph{enhanced flags}. An enhanced flag is a reduction to $N_B=R_u(B)$; forgetting it is a torsor under $T=B/N_B$. With $s$ nodes, the special stack is
\[
 \bigl[(\cA_1^+\times\cA_2^+)/T^s\bigr],
\]
where the torus changes the two branch enhancements at each node simultaneously. All node identifications must be compatible with the same parameter tower. One second component suffices: the extra cycles in the dual graph give gluing choices and impose no further relation between the nodes.

After nearby cycles, pull back to the product and restrict to a fiber over the second factor. This gives an eigencomplex on the first enhanced stack. To pass to ordinary flags, we prove that its monodromies along the enhancement torus are unipotent. The local central-kernel trace at a puncture identifies representation characters of that torus monodromy with traces of the parameter's boundary monodromy. The latter traces are $\dim V$, for every $V$, under \eqref{intro:monodromy}. Characters therefore force all enhancement eigencharacters to be trivial. Compactly supported cohomology along the torus is then nonzero, so forgetting the enhancements preserves a nonzero eigencomplex. This is the step where unipotence, rather than regular unipotence, is sufficient.

Two further arguments control the sheaves throughout. A local calculation with two Hecke modifications detects covectors outside the nilpotent cone. The same calculation, applied to cuts over a trait, proves that nearby cycles cannot vanish when the closure of the nonzero generic stalks meets the special fiber. Separately, the Betti spectral action makes fixed-point Hecke functors exact near a dense parameter: an equivariant frame near its free closed orbit identifies these functors with finite sums of the identity. Moving-leg local constancy then permits perverse cohomology while retaining every tensor and collision comparison.

To reach the given curve, lift it with its marked points over the Witt ring $W(k)$. The finite quotients of $\rho$ extend over root stacks at the marked sections, whose orders are powers of $\ell$. Proper henselian lifting gives a compatible parameter on the geometric generic curve with the same dense image and boundary monodromies. Apply the characteristic-zero construction to that parameter and take geometric nearby cycles on the ordinary Borel-level bundle stack. Proper bounded Hecke modifications make the closure of the nonzero generic-stalk locus meet the special fiber, allowing the detection criterion to prove nonvanishing once more.

These constructions extend the one-point argument of \cite{CT}. The additional work is to prove the detection and perverse statements for the several-level stacks, to keep the other levels transported in the local central-kernel comparison, and to glue one compatible tower at all nodes. The detection criterion and the torus descent are stated separately from the final construction, so their hypotheses and their use in both specializations remain explicit. The nearby-cycle and unramified existence inputs are recalled with their precise scope; the several-point extensions are proved here.

Section~\ref{found:section} fixes the sheaf and Hecke conventions. Section~\ref{geom:section} proves the cotangent and transverse-modification facts used in the detection theorems of Section~\ref{det:section}. Sections~\ref{perv:section} and~\ref{desc:section} establish perverse cohomology and descent from enhanced flags. Section~\ref{con:section} carries out the characteristic-zero construction, and Section~\ref{lift:section} specializes it to prove Theorem~\ref{intro:main}.
